\begin{center}
  \fontsize{14}{17}\selectfont\textbf{ABSTRAK}
\end{center}

\addcontentsline{toc}{chapter}{ABSTRAK}

\vspace{2ex}

\begin{center}
  \fontsize{12}{15}\selectfont\textbf{\tatitle{}}
\end{center}

\vspace{2ex}

\noindent
{\bfseries
\begin{tabular}{@{}p{4.3cm}l}
  Nama Mahasiswa / NRP & : \name{} / \nrp{} \\
  Program Studi           & : \department{} \facultyshort{} -- ITS \\
  Dosen Pembimbing     & : \advisor{} \\
  Dosen Ko-pembimbing  & : \coadvisor{} \\
\end{tabular}
}

\vspace{2ex}

\noindent\textbf{Abstrak}

% Ubah paragraf berikut dengan abstrak dari tugas akhir
Tekanan air pada pompa air rumah tangga cenderung tidak stabil, terutama
ketika beberapa keran dibuka secara bersamaan, sehingga diperlukan suatu
sistem yang mampu menjaga tekanan air tetap stabil. Solusi yang umum
digunakan berupa alat tambahan seperti
\emph{pressure switch} yang bekerja pada sisi keluaran pompa. Penelitian ini
menawarkan pendekatan berbeda, yaitu mengendalikan tekanan dari sisi suplai
daya melalui rangkaian PWM AC \emph{chopper} topologi diode \emph{bridge}
dengan pengendali berbasis \emph{neural network}. Pengendali yang digunakan
adalah \emph{Multi-Layer Perceptron} (MLP) 3-4-1 dengan masukan galat
tekanan, perubahan galat, dan \emph{duty cycle}, serta keluaran berupa
perubahan \emph{duty cycle}, yang dilatih dari data operasi \emph{closed-loop}
(\emph{behavior cloning}). Hasil pengujian menunjukkan rangkaian PWM AC \emph{chopper} mampu mengatur
tegangan RMS keluaran secara linear terhadap \emph{duty cycle} dengan selisih
di bawah 1\% terhadap nilai teoritis, dengan \emph{dead time} terukur
1,48--1,49~$\mu$s sesuai rancangan. Model \emph{neural network} mencapai
koefisien determinasi $R^2 = 0{,}915$ pada data uji, menunjukkan akurasi
prediksi aksi kendali yang baik. Pada pengujian sistem, pengendali mampu menjaga tekanan pada \emph{setpoint}
dengan galat tunak $\leq 0{,}009$~bar untuk variasi satu sampai empat keran
terbuka. Pada pengujian perubahan \emph{setpoint} (0,25--0,35~bar),
\emph{settling time} berkisar 4--7~detik, sedangkan pada pengujian rejeksi
gangguan perubahan beban keran, tekanan pulih dalam 9--15~detik. Keterbatasan
terjadi saat seluruh keran tertutup dan saat gangguan pembukaan empat keran
serentak, di mana tekanan minimum pompa melampaui \emph{setpoint} sehingga
aktuator mengalami saturasi fisik. Dengan demikian, sistem pengendalian tekanan
pompa air berbasis \emph{neural network} dari sisi suplai terbukti bekerja
terhadap variasi beban.

\vspace{1ex}
\noindent\textbf{Kata kunci: \emph{PWM AC Chopper, Neural Network, Tekanan Pompa Air, Multi-Layer Perceptron}.}
